% ------------------------------------------------------------------
% 1_intro.tex — body-only (\input into root.tex)
% ------------------------------------------------------------------
\section{Introduction}
\label{sec:intro}

Robots and autonomous vehicles must perceive their surroundings under night, fog, rain, and sensor failure---conditions in which no single sensor is dependable. Multimodal suites that pair cameras with depth, LiDAR, and event sensors offer complementary evidence, and benchmarks such as DELIVER~\cite{cmnext2023} and MUSES~\cite{muses2024} have made adverse-condition multimodal segmentation a measurable problem. The central difficulty is not only \emph{how} to fuse heterogeneous inputs but \emph{how much to trust each sensor at each location}: sensor competence varies per condition, per region, and---critically---per semantic class, so a fusion policy that is fixed or merely scene-global discards exactly the complementarity that motivates the sensor suite.

Current state-of-the-art systems answer this question with additional supervision and bespoke fusion backbones. DGFusion~\cite{dgfusion2026}, the strongest published method on the DELIVER test split, derives its reliability signal from ground-truth log-depth supervision combined with condition meta-text through a language--vision encoder; CAFuser~\cite{cafuser2025} requires condition annotations and CLIP-based~\cite{clip2021} contrastive training; HyperDUM~\cite{hyperdum2025} builds label-derived prototypes. Orthogonally, the dominant architecture line---CMX~\cite{cmx2023}, CMNeXt~\cite{cmnext2023}, StitchFusion~\cite{stitchfusion2024}, GeminiFusion~\cite{geminifusion2024}---trains full multimodal encoders end-to-end, forgoing the robustness of large frozen vision foundation models (VFMs) whose self-supervised features transfer strongly to dense prediction~\cite{dinov3_2025}.

We take the opposite route (Fig.~\ref{fig:teaser}). \emph{ReliaDINO} encodes all four modalities with a \emph{single frozen} DINOv3 ViT-L/16~\cite{dinov3_2025}; modality identity enters only through independent per-modality low-rank adapters~\cite{lora2022} on the attention projections ($\approx$3.1M parameters in total). Light per-modality decoders yield calibrated class posteriors from which all reliability signals are computed---without any supervision beyond the segmentation labels. Two shared cross-modal attention layers exchange information across modalities, and reliability is applied where class decisions are formed: a calibrated competence gate with a training-free veto floor fuses the modality features, and a per-class reliability-anchored router (PRR) re-weights the per-modality posteriors, letting each class rely on the sensor that actually sees it.

Perhaps our most instructive finding is negative. The seemingly natural injection site for reliability---an additive, pre-softmax bias on cross-modal attention logits, a placement adjacent to recent attention-modulation work~\cite{sam2long2024,sae2026,primed2026}---is consistently ineffective in our setting: across five model generations spanning SAM2-based~\cite{sam2_2024,memorysam2025} and DINOv3-based architectures, such biases changed predictions negligibly or hurt significantly, whereas the identical signals contribute measurably at the output gate and router. Our analysis attributes this to per-modality low-rank adaptation itself: once adapters align the modality representations, attention no longer routes pathologically, and suppressing unreliable keys becomes a no-op---raising \emph{where to inject reliability} as a design question in its own right.

Our contributions are:
\begin{itemize}
  \item \textbf{C1 --- Frozen-VFM modality multiplexing.} A single frozen DINOv3 ViT-L/16 serves all four modalities via per-modality LoRA; to our knowledge the first DINOv3-based multimodal segmentation model, with 46.8M trainable parameters (13.4\%)---fewer than the total of a CMNeXt-B2 baseline.
  \item \textbf{C2 --- Reliability at the decision level.} Calibrated competence-gated fusion (+0.26 test mIoU at full resolution) and a per-class reliability-anchored router; reliability acts at the output gate and router rather than as an attention-logit bias.
  \item \textbf{C3 --- Supervision minimality.} Segmentation labels only; no ground-truth depth, condition annotations, or text encoders, in contrast to the methods we surpass.
  \item \textbf{C4 --- Analysis.} A five-generation controlled negative result on attention-logit reliability biasing; the first per-class audit of the DELIVER \emph{test} split, exposing a benchmark ceiling (Wall, Bridge, and Other are near-zero even for the official baseline \todo{confirm official CMNeXt per-class test numbers, M0-a}); and a strictly validation-based checkpoint-selection protocol.
\end{itemize}

On DELIVER with all four modalities, ReliaDINO attains \todo{final test mIoU; best so far 57.14} on the test split (previous best: DGFusion, 56.71) and \todo{final val mIoU; best so far 67.74} on the validation split, while its no-router variant reaches 68.19 val / 57.60 test mIoU---evidence that a frozen foundation model, minimally adapted and fused by self-derived reliability, outperforms specialized, extra-supervised fusion architectures.
